% Source: https://artofproblemsolving.com/community/c3h1867617_a_diophantine_equation
% Stable identity: aops-topic-1867617-q01
The equation is $\dfrac{1}{a}+\dfrac{1}{b}+\dfrac{1}{c}=\dfrac{1}{2}$. Since $a\le b\le c$, we have $\dfrac{1}{2}\le\dfrac{3}{a}$, so $a\le6$; and $\dfrac{1}{a}<\dfrac{1}{2}$, so $a\ge3$.

For fixed $a$, the equation $\dfrac{1}{b}+\dfrac{1}{c}=\dfrac{1}{m}$ (with $m$ a positive integer) is equivalent to $(b-m)(c-m)=m^2$ with $b,c>m$. We use this in each case.
\begin{itemize}
\item $a=3$: $\dfrac{1}{b}+\dfrac{1}{c}=\dfrac{1}{6}$, so $(b-6)(c-6)=36$. With $b\le c$: $b-6\in\{1,2,3,4,6\}$, giving $(b,c)=(7,42),(8,24),(9,18),(10,15),(12,12)$; all have $b\ge3$.
\item $a=4$: $\dfrac{1}{b}+\dfrac{1}{c}=\dfrac{1}{4}$, so $(b-4)(c-4)=16$, giving $(b,c)=(5,20),(6,12),(8,8)$.
\item $a=5$: $\dfrac{1}{b}+\dfrac{1}{c}=\dfrac{3}{10}$. With $5\le b\le c$, $\dfrac{2}{b}\ge\dfrac{3}{10}$ gives $b\le6$. $b=5$ gives $\dfrac{1}{c}=\dfrac{1}{10}$, $c=10$; $b=6$ gives $\dfrac{1}{c}=\dfrac{2}{15}$, not a unit fraction.
\item $a=6$: $\dfrac{1}{b}+\dfrac{1}{c}=\dfrac{1}{3}$ with $6\le b\le c$ forces $b=c=6$.
\end{itemize}
So there are $10$ triples:
\[
\boxed{(3,7,42),(3,8,24),(3,9,18),(3,10,15),(3,12,12),(4,5,20),(4,6,12),(4,8,8),(5,5,10),(6,6,6)}.
\]
